\chapter{组合与计数}\label{ch:g12:comb}

\omterm{def:g12:comb:combination}{组合}数学是不借助列举而计数的艺术。它的两条基本原理——
把互不相交的备选方案的个数相加，把独立选择的个数相乘——
足以计数有限集的\omterm{def:g12:comb:tuples}{排列}、\omterm{def:g12:comb:permutation}{全排列}与子集，
并最终导出二项式定理。

\section{两个计数原理}

记 $\abs{E}$ 为有限集
$E$ 的元素个数（即\emph{基数}）。

\begin{proposition}[加法原理]\label{prop:g12:comb:addition}
若有限集 $E$ 被划分为子集 $A_1, \dots, A_k$（两两
不交，且\omterm{def:g10:numbers:interunion}{并集}为 $E$），则
\[
\abs{E} = \abs{A_1} + \abs{A_2} + \dots + \abs{A_k}.
\]
\end{proposition}

\begin{proposition}[乘法原理]\label{prop:g12:comb:multiplication}
若一个对象由 $k$ 次\omterm{def:g12:limcont:continuity}{连续}选择构成，第一次选择有 $n_1$ 种选法，
且\emph{无论前面的选择如何}，$n_i$ 为
第 $i$ 次选择的选法数，则可构成的对象个数
为 $n_1 \times n_2 \times \dots \times n_k$。
\end{proposition}

\begin{proof}
两条原理均可对 $k$ 用归纳法证明；第二条中 $k = 2$ 的情形
相当于按行数出一个矩形阵列的元素个数。
\end{proof}

\begin{example}
某餐厅提供 4 种开胃菜、6 种主菜、3 种甜点：$4 \times 6 \times 3 =
72$ 种不同的三道式套餐。
\end{example}

\section{元组、全排列与阶乘}

\begin{definition}[$k$-元组]\label{def:g12:comb:tuples}
一个\emph{$k$-元组}是集合 $E$ 的有序列表 $(x_1, \dots, x_k)$，
其元素均取自 $E$，允许重复。若 $k$-元组的元素\emph{互不相同}，
则称为\emph{排列}\index{排列}，即 $k$ 个取自 $E$ 的元素。
\end{definition}

\begin{proposition}\label{prop:g12:comb:tuples}
设 $\abs E = n$。$k$-元组（元素取自 $E$）共有 $n^k$ 个。
$k$ 个元素（取自 $E$）的\omterm{def:g12:comb:tuples}{排列}数（$0 \leq k \leq n$）为
\[
n(n-1)(n-2)\cdots(n-k+1) = \frac{n!}{(n-k)!},
\]
其中 $n! = 1 \times 2 \times \dots \times n$（且 $0! = 1$）称为
$n$ 的\emph{阶乘}\index{阶乘}。
\end{proposition}

\begin{proof}
由乘法原理，对 $k$-元组，每步有 $n$ 种选法，
共 $k$ 步；对\omterm{def:g12:comb:tuples}{排列}，有 $n$ 种选法选取 $x_1$，再有 $n - 1$ 种
选取 $x_2$（已用去一个元素），……，$n - k + 1$ 种选取 $x_k$。
\end{proof}

\begin{definition}[全排列]\label{def:g12:comb:permutation}
\emph{全排列}\index{全排列}是指 $E$ 的全部 $n$ 个元素的一个\omterm{def:g12:comb:tuples}{排列}，
即 $E$ 的一个排序，换言之 $E$ 的元素按某一次序排成一列。
由\cref{prop:g12:comb:tuples}（取 $k = n$）可知，
$n$ 元集合的全\omterm{def:g12:comb:tuples}{排列}个数为 $n!$。
\end{definition}

\begin{example}
五名选手跑完一场比赛，共有 $5! = 120$ 种不同的名次。
可能的领奖台（前三名）个数为 $5 \times 4 \times 3 = 60$。
\end{example}

\section{组合与二项式系数}

\begin{definition}[组合]\label{def:g12:comb:combination}
\emph{组合}\index{组合}是指 $k$ 个取自 $E$ 的元素（不计顺序，不重复）构成的子集，
即 $E$ 中含 $k$ 个元素的子集，其个数记为
$\dbinom{n}{k}$\index{二项式系数}，读作“$n$ 选 $k$”。
\end{definition}

\begin{theorem}\label{thm:g12:comb:binomformula}
对 $0 \leq k \leq n$，有：
\[
\binom{n}{k} = \frac{n!}{k!\,(n-k)!} .
\]
\end{theorem}

\begin{proof}
以两种方式计算 $k$ 个元素（取自 $E$）的\omterm{def:g12:comb:tuples}{排列}数。一是直接计算：
$\frac{n!}{(n-k)!}$。另一种方法是先选定由这些元素构成的子集，
有 $\binom nk$ 种选法，再排序有 $k!$ 种，由乘法原理
得 $\binom{n}{k}\,k!$。两式相等，
即 $\binom nk = \frac{n!}{k!(n-k)!}$。
\end{proof}

\begin{proposition}[基本恒等式]\label{prop:g12:comb:identities}
对 $0 \leq k \leq n$，有：
\[
\binom{n}{0} = \binom{n}{n} = 1, \qquad
\binom{n}{1} = n, \qquad
\binom{n}{k} = \binom{n}{n-k},
\]
以及\emph{帕斯卡法则}\index{帕斯卡法则}：对 $1 \leq k \leq n-1$，
\[
\binom{n}{k} = \binom{n-1}{k-1} + \binom{n-1}{k}.
\]
\end{proposition}

\begin{proof}
对称性 $\binom nk = \binom{n}{n-k}$ 成立，因为取补集
使 $k$ 元子集与 $(n-k)$ 元子集一一对应。对于
\omterm{prop:g12:comb:identities}{帕斯卡法则}，固定元素 $a \in E$，将 $k$ 元子集分为两类：
含 $a$ 的一类——它们由 $a$ 添入一个 $(k-1)$ 元子集得到，
$E \setminus \{a\}$ 的子集，共 $\binom{n-1}{k-1}$ 个——以及
不含 $a$ 的一类，即 $k$ 元子集（取自 $E \setminus \{a\}$），
它们共 $\binom{n-1}{k}$ 个；由加法原理
即得。
\end{proof}

\omterm{prop:g12:comb:identities}{帕斯卡法则}逐行生成各项系数——\emph{帕斯卡
三角形}：每一项是其上方两项之和。

\begin{omfigure}
\begin{tikzpicture}[x=0.95cm, y=0.8cm, font=\small]
  \node at (0,0) {$1$};
  \node at (-0.5,-1) {$1$};   \node at (0.5,-1) {$1$};
  \node at (-1,-2) {$1$};     \node at (0,-2) {$2$};
  \node at (1,-2) {$1$};
  \node at (-1.5,-3) {$1$};   \node at (-0.5,-3) {$3$};
  \node at (0.5,-3) {$3$};    \node at (1.5,-3) {$1$};
  \node at (-2,-4) {$1$};
  \node[circle, draw=omProp, inner sep=1.5pt] (p1) at (-1,-4) {$4$};
  \node[circle, draw=omProp, inner sep=1.5pt] (p2) at (0,-4) {$6$};
  \node at (1,-4) {$4$};      \node at (2,-4) {$1$};
  \node at (-2.5,-5) {$1$};   \node at (-1.5,-5) {$5$};
  \node[circle, draw=omThm, inner sep=1.5pt] (s) at (-0.5,-5) {$10$};
  \node at (0.5,-5) {$10$};   \node at (1.5,-5) {$5$};
  \node at (2.5,-5) {$1$};
  \draw[-Stealth, omProp] (p1) -- (s);
  \draw[-Stealth, omProp] (p2) -- (s);
\end{tikzpicture}

{\small 帕斯卡三角形，第 $n = 0$ 至 $5$ 行：\omterm{prop:g12:comb:identities}{帕斯卡法则}
$\binom{4}{1} + \binom{4}{2} = \binom{5}{2}$ 的体现。}
\end{omfigure}

\begin{theorem}[二项式定理]\label{thm:g12:comb:binomial}
\index{二项式定理}
对所有 $a, b \in \R$（或 $\C$）及 $n \in \N$：
\[
(a+b)^n = \sum_{k=0}^{n} \binom{n}{k}\, a^{k}\, b^{\,n-k} .
\]
\end{theorem}

\begin{proof}
将乘积 $(a+b)(a+b)\cdots(a+b)$（共 $n$ 个因式）\omterm{def:g10:algebra:expand}{展开}，
\omterm{def:g10:algebra:expand}{展开}式的每一项在每个因式中取 $a$ 或 $b$，得到 $a^k b^{n-k}$，
其中 $k$ 为取 $a$ 的因式个数。
这 $k$ 个因式在 $n$ 个因式中的选法数为 $\binom nk$，
故 $a^k b^{n-k}$ 的系数即为此选法数。
\end{proof}

\begin{corollary}\label{cor:g12:comb:sums}
$\displaystyle\sum_{k=0}^{n} \binom{n}{k} = 2^n$
\quad 且 \quad
$\displaystyle\sum_{k=0}^{n} (-1)^k\binom{n}{k} = 0$ \ ($n \geq 1$)。
\end{corollary}

\begin{proof}
在二项式定理中取 $a = b = 1$，再取 $a = -1$，$b = 1$。第一个
恒等式还有直接的含义：$n$ 元集合共有 $2^n$ 个子集
（每个元素取或不取，由乘法原理得），并按子集大小分类。
\end{proof}

\begin{method}[选取正确的模型]\label{met:g12:comb:model}
计数之前，先回答两个问题：\emph{顺序是否重要？}以及
\emph{是否允许重复？}
\begin{center}
\begin{tabular}{l|cc}
 & 顺序有关 & 顺序无关 \\
\hline
允许重复 & $n^k$（元组） & （大学）\\
不允许重复 & $\frac{n!}{(n-k)!}$（\omterm{def:g12:comb:tuples}{排列}） & $\binom nk$（子集）
\end{tabular}
\end{center}
从缸中摸球：\emph{有放回、按顺序} $\to$ 元组；
\emph{无放回、按顺序} $\to$ \omterm{def:g12:comb:tuples}{排列}；\emph{一把全部
取出} $\to$ 子集。
\end{method}

\section{练习}

\begin{exercise}[$\star$]\label{exo:g12:comb:1}
某车牌由 2 个字母（A--Z）、3 个数字和 2 个字母依次组成。
共有多少种可能的车牌？其中没有重复字符的有多少种？
\end{exercise}

\begin{exercise}[$\star$]\label{exo:g12:comb:2}
计算 $\dbinom{8}{3}$、$\dbinom{10}{8}$，并化简
$\dfrac{\binom{n}{2}}{\binom{n+1}{2}}$。
\end{exercise}

\begin{exercise}[$\star$]\label{exo:g12:comb:3}
某班有 30 名学生，需选出 4 名学生组成委员会，再在委员会中选出
一名主席和一名财务负责人（同一人不能
兼任两个职务）。共有多少种可能的结果？
\end{exercise}

\begin{exercise}[$\star$]\label{exo:g12:comb:4}
利用二项式定理\omterm{def:g10:algebra:expand}{展开} $(x + 2)^5$ 和 $(1 - x)^6$。
求 $x^3$ 在 $(2x + 3)^7$ 中的系数。
\end{exercise}

\begin{exercise}[$\star\star$]\label{exo:g12:comb:5}
标准扑克牌局中的一手牌由 52 张牌中的 5 张组成。
\begin{enumerate}
  \item 共有多少种可能的手牌？
  \item 恰含一张 A 的手牌有多少种？至少含一张 A 的呢？
  \item “葫芦”手牌（三张同点数、两张
        另一点数）有多少种？
\end{enumerate}
\end{exercise}

\begin{exercise}[$\star\star$]\label{exo:g12:comb:6}
求单词 MATH 的字母重排（即字母的重新\omterm{def:g12:comb:tuples}{排列}，不论是否有意义）
共有多少种？单词 BANANA 呢？
（提示：对 BANANA，先安排三个 A 的位置。）
\end{exercise}

\begin{exercise}[$\star\star$]\label{exo:g12:comb:7}
用两种方法证明恒等式 $k\dbinom{n}{k} = n\dbinom{n-1}{k-1}$
（$1 \leq k \leq n$），一是用\omterm{prop:g12:comb:tuples}{阶乘}公式，二是对
“$k$ 人委员会及其主席”这类数对（从 $n$ 人中选取）作两种计数。
\end{exercise}

\begin{exercise}[$\star\star$]\label{exo:g12:comb:8}
平面上一条路径从 $(0,0)$ 出发，每次沿正东或正北走单位步长，
到达 $(m, n)$。证明此类路径共有 $\dbinom{m+n}{m}$ 条。
\end{exercise}

\begin{exercise}[$\star\star\star$]\label{exo:g12:comb:9}
证明\emph{范德蒙德恒等式}：对 $0 \leq k \leq m + n$，
\[
\binom{m+n}{k} = \sum_{j=0}^{k} \binom{m}{j}\binom{n}{k-j},
\]
通过计数一个集合的 $k$ 元子集来证明，该集合分成含 $m$ 个元素的一组与
含 $n$ 个元素的一组。由此推出
$\displaystyle\sum_{j=0}^{n}\binom{n}{j}^{\!2} = \binom{2n}{n}$.
\end{exercise}

\begin{exercise}[$\star\star\star$]\label{exo:g12:comb:10}
利用二项式定理证明：对所有 $n \geq 1$，
\[
\sum_{k=1}^{n} k \binom{n}{k} = n\,2^{n-1}.
\]
（提示：可对 $(1+x)^n$ 求导，也可利用\cref{exo:g12:comb:7}。）
\end{exercise}

\section{问题：两次计数的艺术}

\begin{problem}[{周末问题——隔板法、错排帽子，
以及用两种方式计数同一事物所证明的恒等式}]
\label{pb:g12:comb:1}
\omterm{def:g12:comb:combination}{组合}计数中最深刻的技巧，其实朴素得令人意外：对同一集合计数，
用两种不同的方法各计数一次，再令两个答案相等。
本题用到\cref{met:g12:comb:model}所列的模型，
并增加一项本章课程未用到的技巧，
即冰淇淋计数中的\emph{隔板法}，
然后精确计数著名的\emph{错排帽子}问题，
并在帽子堆底部找到数 $\frac1\eu$，
这是它在本书中第三次出现。

\medskip
\noindent\textbf{第一部分——选择模型。}
\begin{enumerate}
  \item 计数由 $2$ 个字母后接
        $3$ 个数字组成的车牌数；再计数 BANANA 的字母重排数。
  \item 从 $32$ 张牌的牌组中计数 $5$ 张牌的手牌数；再计数
        恰含 $2$ 张 A（共 $4$ 张 A）的手牌数。
  \item 机器人每次只能向右或向上走一个单位，从 $(0,0)$ 走到
        $(4,3)$，共有多少条路径？（将路径编码为由 R 与 U 组成的
        字符串。）
  \item 由二项式定理（\cref{thm:g12:comb:binomial}）\omterm{def:g10:algebra:expand}{展开} $(1 + x)^4$，
        再分别代入 $x = 1$ 与 $x = -1$，
        可得出关于
        $\binom nk$ 的哪两个恒等式？
  \item 用双重计数法证明
        $k\binom nk = n\binom{n-1}{k-1}$（以两种方式计数带主席的委员会）
        ，并由此推出
        $\sum_{k=0}^{n} k\binom nk = n\,2^{n-1}$.
\end{enumerate}

\medskip
\noindent\textbf{第二部分——隔板法。}
\begin{enumerate}[resume]
  \item 一家冰淇淋店售卖 $4$ 种口味，你点了 $10$
        勺（口味可以重复，杯中顺序无关）。将一种点法编码为
        由 $10$ 个星号（勺）与 $3$ 个隔板（口味变换）
        组成的一排，
        并计数这样的点法。
  \item 计数满足
        $x + y + z = 12$ 的非负\omterm{def:g10:numbers:sets}{整数}三元组个数。
  \item 计数由\emph{正}\omterm{def:g10:numbers:sets}{整数}构成的三元组，
        满足 $x + y + z = 12$（令 $x = 1 + x'$ 等作代换）。
  \item 在
        $(a + b + c)^5$ 的\omterm{def:g10:algebra:expand}{展开}式中共有多少个不同的单项式？
  \item 检验该方法：用公式计数点 $3$ 勺冰淇淋的
        点法数（口味共 $2$ 种），然后将它们全部列出，
        并与之比较。
  \item 准确指出“勺子相同”是在哪一步进入编码的——
        并计数：若勺子按顺序食用（位置互异），
        结果会有什么不同？
        请用 \cref{met:g12:comb:model} 的检查清单核对。
\end{enumerate}

\medskip
\noindent\textbf{第三部分——错排帽子。}
\emph{错排}是指将 $n$ 顶帽子重新分配给其
$n$ 位主人，且\emph{无人}拿到自己的帽子；记
$D_n$ 为其个数。（\Cref{pb:g11:prob:1} 指出，平均有一位宾客
能拿回自己的帽子——现在我们精确计数完全
倒霉的情形。）
\begin{enumerate}[resume]
  \item 通过列举计算 $D_1$、$D_2$、$D_3$，
        并（耐心或巧妙地）计算 $D_4$。
  \item 证明递推关系
        $D_n = (n - 1)\left(D_{n-1} + D_{n-2}\right)$：宾客 1
        拿到某顶帽子 $k \neq 1$（有 $n - 1$ 种选择）；
        再按宾客 $k$ 是否拿到 1 号帽子分类讨论。
        检验它能否给出 $D_4$，并计算 $D_5$。
  \item 对 $n = 3$，用容斥原理（减去至少固定一顶帽子的分配数，
        再加回重复计数的部分），
        证明
        $D_3 = 3!\left(1 - \frac{1}{1!} + \frac{1}{2!} -
        \frac{1}{3!}\right)$，并给出一般公式。
  \item 计算 $\frac{D_5}{5!}$，并与
        $\frac1\eu \approx 0.3679$ 作比较：随机洗牌后，
        一场大型聚会中帽子被完全错排的\omterm{def:g10:proba:distribution}{概率}为
        $\frac1\eu$——这是该常数的第三次登场，
        继彩票问题以及 \cref{pb:g12:exp:1} 中的秘书问题之后。
        （原因：第 14 题中的公式是一个著名级数的开端，
        该级数表示 $\eu^{-1}$，
        将在大学分册中讲述。）
  \item 在 $10$ 位朋友之间进行“秘密圣诞老人”抽签：名字
        均匀随机抽取。抽签有效（即无人抽到自己）的\omterm{def:g10:proba:distribution}{概率}是多少，
        以及该小组平均应预计
        需要重抽多少次？
\end{enumerate}

\medskip
\noindent\textbf{第四部分——两次计数，两次制胜。}
\begin{enumerate}[resume]
  \item 握手引理：在任一聚会中，对所有宾客求和，
        即每位宾客握手的次数之和，每次握手恰好被计入两次。由此推出
        \emph{握手次数为奇数的宾客人数
        总是偶数}——
        并在三人聚会中检验该结论是否合理。
  \item 用归纳法证明明珠恒等式
        $1^3 + 2^3 + \dots + n^3 = (1 + 2 + \dots + n)^2$，
        并验证 $n = 3$ 的情形。（小高斯的求和式，
        平方后即得立方和。）
  \item 借助路径证明范德蒙德恒等式（\cref{exo:g12:comb:9}）：
        将 $\binom{2n}{n}$ 解释为第 3 题那类、从 $(0,0)$ 到 $(n,n)$ 的格点路径，
        将每条路径在反对角线上的交点处截开，
        并说明
        $\sum_j \binom nj^2$ 如何出现。
  \item 总结——计数的四个动作，每个都配本题中的例子：
        分步相乘，分类相加；
        巧妙编码（隔板法、路径字符串）；
        对同一事物计数两次（带主席的委员会、握手）；
        减去不需要的部分，修正
        重复计数（错排）。注意计数接下来还将用于
        \omterm{def:g10:proba:distribution}{概率}，以及
        “矩阵与图”一章中的路径。
\end{enumerate}
\end{problem}
